\documentclass[12pt,a4paper]{article}

% Сообщение Ярославу:
% В работе хорошо выполнены задания за второе октября и два последних задания за третье октября.
% В первом задании за третье октября в итоговый ответ попала лишняя серия значений,
% при которых котангенс не определён. Задания за четвёртое октября в присланных
% страницах решения отсутствуют, поэтому они разобраны как неполные.

\usepackage[a4paper,margin=2cm]{geometry}
\usepackage[T2A]{fontenc}
\usepackage[utf8]{inputenc}
\usepackage[russian]{babel}
\usepackage{amsmath,amssymb,mathtools}
\usepackage{array,booktabs,longtable}
\usepackage{xcolor}
\usepackage{hyperref}
\usepackage{tikz}
\usepackage{fancyhdr}

\hypersetup{
  colorlinks=true,
  linkcolor=blue!55!black,
  urlcolor=blue!55!black
}

\setlength{\parindent}{0pt}
\setlength{\parskip}{0.55em}
\renewcommand{\arraystretch}{1.25}

\pagestyle{fancy}
\fancyhf{}
\fancyfoot[C]{\small Лёвин Артём Александрович эксклюзивно для Ярослава Гаврилова}
\renewcommand{\headrulewidth}{0pt}
\renewcommand{\footrulewidth}{0pt}

\newcommand{\err}{\textcolor{red}{Ошибка:}}
\newcommand{\idea}{\textcolor{blue!60!black}{Ключевая идея:}}
\newcommand{\result}{\textcolor{teal!60!black}{Итог:}}

\begin{document}

\thispagestyle{empty}
\vspace*{\fill}
\begin{center}
{\LARGE\bfseries Ревью домашней работы\par}
\vspace{1.2cm}
{\large Автор: Лёвин Артём Александрович\par}
\vspace{0.7cm}
{\large 3 октября 2026 года\par}
\end{center}
\vspace*{\fill}
\begin{center}
\begin{tikzpicture}[scale=1]
  \draw[blue!35!black, line width=0.7pt]
  (0,0) .. controls (2,0.7) and (4,-0.7) .. (6,0);
\end{tikzpicture}
\end{center}
\newpage

{\small
\setlength{\tabcolsep}{3pt}
\begin{longtable}{>{\raggedright\arraybackslash}p{0.11\linewidth}
                  >{\raggedright\arraybackslash}p{0.15\linewidth}
                  >{\raggedright\arraybackslash}p{0.28\linewidth}
                  >{\raggedright\arraybackslash}p{0.18\linewidth}
                  >{\raggedright\arraybackslash}p{0.18\linewidth}}
\toprule
№ задания & Тема & Суть ошибки & Ваш ответ & Правильный ответ \\
\midrule
\endfirsthead
\toprule
№ задания & Тема & Суть ошибки & Ваш ответ & Правильный ответ \\
\midrule
\endhead
\hyperref[task:0310-1]{03.10, №1} &
Произведение тригонометрических множителей с радикалом &
В ответ включены значения \(x=\pi n\), но при них \(\operatorname{ctg}x\) не определён. &
\(\frac{\pi}{2}+2\pi n;\ \frac{\pi}{4}+2\pi n;\ \pi n\) &
\(\frac{\pi}{2}+2\pi n;\ \frac{\pi}{4}+2\pi n\) \\
\midrule
\hyperref[task:0410-1]{04.10, №1} &
Тригоно\-метрическая дробь &
Решение и ответ не представлены. &
не указан &
решений нет \\
\midrule
\hyperref[task:0410-2]{04.10, №2} &
Дробь с квадратным корнем &
Решение и ответ не представлены. &
не указан &
\(x=\pi n,\ n\in\mathbb Z\) \\
\midrule
\hyperref[task:0410-3]{04.10, №3} &
Произведение с тангенсом и радикалом &
Решение и ответ не представлены. &
не указан &
\(x=2\pi n\) или \(x=\frac{\pi}{4}+2\pi n,\ n\in\mathbb Z\) \\
\bottomrule
\end{longtable}
}

\newpage
\thispagestyle{empty}
\vspace*{\fill}
\begin{center}
{\LARGE\bfseries Разбор ошибок}
\end{center}
\vspace*{\fill}
\newpage

\phantomsection\label{task:0310-1}
\section*{03.10.2026 — задание №1}

\textbf{Текст задания}

Решить уравнение
\[
(\sin x-1)(\operatorname{ctg}x-1)\sqrt{\sin x}=0.
\]

\textbf{Ваше решение}

В решении были отдельно рассмотрены множители:
\[
\sin x-1=0,\qquad
\operatorname{ctg}x-1=0,\qquad
\sqrt{\sin x}=0.
\]
Для первого множителя получено
\[
x=\frac{\pi}{2}+2\pi n.
\]
Для второго сначала записано
\[
x=\frac{\pi}{4}+\pi n,
\]
а в итоговом ответе оставлена серия
\[
x=\frac{\pi}{4}+2\pi n.
\]
Для третьего множителя получено
\[
\sin x=0,\qquad x=\pi n.
\]

\textbf{Ваш ответ}

\[
x=\frac{\pi}{2}+2\pi n,\qquad
x=\frac{\pi}{4}+2\pi n,\qquad
x=\pi n,\qquad n\in\mathbb Z.
\]

\textbf{\err\ в итоговый ответ включена серия, не принадлежащая области допустимых значений.}

\textbf{Где именно возникает ошибка}

Последний корректный шаг --- решение уравнения
\[
\sqrt{\sin x}=0,
\]
из которого формально следует
\[
\sin x=0,\qquad x=\pi n.
\]
Первый неверный шаг --- перенос этих значений в окончательный ответ без проверки исходного выражения.

\textbf{Почему этот шаг неверен}

В исходном уравнении одновременно присутствуют квадратный корень
\(\sqrt{\sin x}\) и котангенс \(\operatorname{ctg}x\).

Для существования квадратного корня необходимо
\[
\sin x\ge 0.
\]
Для существования котангенса необходимо
\[
\sin x\ne 0.
\]
Оба условия должны выполняться одновременно, поэтому область допустимых значений задаётся условием
\[
\sin x>0.
\]

При \(x=\pi n\) имеем
\[
\sin(\pi n)=0.
\]
Следовательно, \(\operatorname{ctg}(\pi n)\) не определён. Значит, исходное выражение при таких значениях \(x\) вообще не имеет смысла, и они не могут быть корнями.

Важно различать два утверждения: отдельный множитель \(\sqrt{\sin x}\) действительно обращается в нуль при \(\sin x=0\), но произведение можно считать равным нулю только в тех точках, где определены \emph{все} множители исходного выражения.

\textbf{\idea} сначала нужно найти общую область допустимых значений всего выражения, а уже затем решать уравнение произведения, пересекая решения каждого множителя с этой областью.

\textbf{Исправление от последнего правильного шага}

Из
\[
\sqrt{\sin x}=0
\]
получаем
\[
x=\pi n.
\]
Но эти значения не удовлетворяют условию
\[
\sin x>0.
\]
Поэтому вся серия \(x=\pi n\) должна быть исключена.

\textbf{Полное правильное решение}

Сначала найдём область допустимых значений.

Из радикала:
\[
\sin x\ge 0.
\]
Из котангенса:
\[
\sin x\ne 0.
\]
Следовательно,
\[
\sin x>0.
\]

Теперь произведение равно нулю, если хотя бы один множитель равен нулю, причём полученные значения принадлежат области допустимых значений.

\textbf{1. Первый множитель}
\[
\sin x-1=0,
\]
\[
\sin x=1,
\]
\[
x=\frac{\pi}{2}+2\pi n,\qquad n\in\mathbb Z.
\]
Здесь \(\sin x=1>0\), поэтому все найденные значения допустимы.

\textbf{2. Второй множитель}
\[
\operatorname{ctg}x-1=0,
\]
\[
\operatorname{ctg}x=1.
\]
Общее решение:
\[
x=\frac{\pi}{4}+\pi n,\qquad n\in\mathbb Z.
\]
Но нужно оставить только те значения, для которых \(\sin x>0\).

При
\[
x=\frac{\pi}{4}+2\pi n
\]
синус положителен, а при добавлении нечётного числа \(\pi\) знак синуса меняется. Поэтому допустимая серия:
\[
x=\frac{\pi}{4}+2\pi n,\qquad n\in\mathbb Z.
\]

\textbf{3. Третий множитель}
\[
\sqrt{\sin x}=0,
\]
\[
\sin x=0,
\]
\[
x=\pi n,\qquad n\in\mathbb Z.
\]
Но эти значения не входят в область допустимых значений, так как при них котангенс не определён. Поэтому эта серия исключается.

\textbf{Проверка}

Для
\[
x=\frac{\pi}{2}+2\pi n
\]
имеем \(\sin x=1\), поэтому первый множитель равен нулю, а котангенс определён.

Для
\[
x=\frac{\pi}{4}+2\pi n
\]
имеем \(\sin x>0\) и \(\operatorname{ctg}x=1\), поэтому второй множитель равен нулю и исходное выражение определено.

Значения \(x=\pi n\) проверку не проходят, потому что при них \(\operatorname{ctg}x\) не существует.

\textbf{\result}
\[
\boxed{
x=\frac{\pi}{2}+2\pi n
\quad\text{или}\quad
x=\frac{\pi}{4}+2\pi n,\qquad n\in\mathbb Z.
}
\]

\textbf{Задача на закрепление}

Решите уравнение
\[
(2\sin x-1)(\operatorname{ctg}x+1)\sqrt{\sin x}=0.
\]

\newpage
\phantomsection\label{task:0410-1}
\section*{04.10.2026 — задание №1}

\textbf{Текст задания}

Решить уравнение
\[
\frac{\sin 4x}{\sin 2x}=2.
\]

\textbf{Ваше решение}

Решение не представлено.

\textbf{Ваш ответ}

Не указан.

\textbf{\err\ задание не выполнено.}

\textbf{Где именно возникает ошибка}

Последний корректный шаг отсутствует, поскольку решение не начато.

Первый неполный шаг --- отсутствует запись области допустимых значений и дальнейшее преобразование исходного уравнения.

\textbf{Почему это существенно}

В этом уравнении нельзя просто сократить дробь, не зафиксировав условие существования знаменателя. Знаменатель не должен обращаться в нуль:
\[
\sin 2x\ne 0.
\]
Именно это условие в конце исключит все формально найденные значения.

\textbf{\idea} при сокращении тригонометрической дроби условие ненулевого знаменателя сохраняется до конца решения и обязательно проверяется после нахождения кандидатов.

\textbf{Исправление}

Начинаем с области допустимых значений:
\[
\sin 2x\ne 0.
\]
Затем используем формулу двойного угла для синуса:
\[
\sin 4x=2\sin 2x\cos 2x.
\]

\textbf{Полное правильное решение}

Исходное уравнение:
\[
\frac{\sin 4x}{\sin 2x}=2.
\]

Область допустимых значений:
\[
\sin 2x\ne 0.
\]

Подставим
\[
\sin 4x=2\sin 2x\cos 2x.
\]
Получаем
\[
\frac{2\sin 2x\cos 2x}{\sin 2x}=2.
\]
Так как по области допустимых значений \(\sin 2x\ne 0\), можно сократить:
\[
2\cos 2x=2.
\]
Делим обе части на \(2\):
\[
\cos 2x=1.
\]
Отсюда
\[
2x=2\pi n,\qquad n\in\mathbb Z,
\]
поэтому
\[
x=\pi n,\qquad n\in\mathbb Z.
\]

Теперь проверяем область допустимых значений:
\[
\sin 2x=\sin(2\pi n)=0.
\]
Все найденные значения запрещены.

\textbf{Проверка}

При любом \(x=\pi n\) знаменатель исходной дроби равен нулю, поэтому исходное выражение не определено.

\textbf{\result} уравнение не имеет решений.

\textbf{Задача на закрепление}

Решите уравнение
\[
\frac{\sin 6x}{\sin 3x}=2.
\]

\newpage
\phantomsection\label{task:0410-2}
\section*{04.10.2026 — задание №2}

\textbf{Текст задания}

Решить уравнение
\[
\frac{\sqrt{\sin x}}{\cos x}=0.
\]

\textbf{Ваше решение}

Решение не представлено.

\textbf{Ваш ответ}

Не указан.

\textbf{\err\ задание не выполнено.}

\textbf{Где именно возникает ошибка}

Последний корректный шаг отсутствует, поскольку решение не начато.

Первый неполный шаг --- не записаны условия существования квадратного корня и знаменателя, а также условие равенства дроби нулю.

\textbf{Почему это существенно}

Дробь равна нулю тогда и только тогда, когда её числитель равен нулю, а знаменатель при этом не равен нулю. Кроме того, квадратный корень существует только при неотрицательном подкоренном выражении.

Значит, одновременно нужно учитывать:
\[
\sin x\ge 0,\qquad \cos x\ne 0.
\]

\textbf{\idea} для уравнения «дробь равна нулю» сначала приравниваем к нулю числитель, но найденные значения обязательно проверяем по знаменателю и по области допустимых значений радикала.

\textbf{Исправление}

Чтобы дробь была равна нулю, требуется
\[
\sqrt{\sin x}=0
\]
при условии
\[
\cos x\ne 0.
\]

\textbf{Полное правильное решение}

Исходное уравнение:
\[
\frac{\sqrt{\sin x}}{\cos x}=0.
\]

Область допустимых значений:
\[
\sin x\ge 0,\qquad \cos x\ne 0.
\]

Приравниваем числитель к нулю:
\[
\sqrt{\sin x}=0.
\]
Квадратный корень равен нулю только тогда, когда подкоренное выражение равно нулю:
\[
\sin x=0.
\]
Отсюда
\[
x=\pi n,\qquad n\in\mathbb Z.
\]

Проверим знаменатель:
\[
\cos(\pi n)=(-1)^n.
\]
Это число всегда равно \(1\) или \(-1\), следовательно, никогда не равно нулю.

Условие \(\sin x\ge 0\) также выполнено, поскольку в найденных точках \(\sin x=0\).

\textbf{Проверка}

При \(x=\pi n\) числитель равен
\[
\sqrt{0}=0,
\]
а знаменатель равен \(1\) или \(-1\). Поэтому дробь действительно равна нулю.

\textbf{\result}
\[
\boxed{x=\pi n,\qquad n\in\mathbb Z.}
\]

\textbf{Задача на закрепление}

Решите уравнение
\[
\frac{\sqrt{\cos x}}{\sin x}=0.
\]

\newpage
\phantomsection\label{task:0410-3}
\section*{04.10.2026 — задание №3}

\textbf{Текст задания}

Решить уравнение
\[
(\cos 2x-1)(\operatorname{tg}x-1)\sqrt{\cos x}=0.
\]

\textbf{Ваше решение}

Решение не представлено.

\textbf{Ваш ответ}

Не указан.

\textbf{\err\ задание не выполнено.}

\textbf{Где именно возникает ошибка}

Последний корректный шаг отсутствует, поскольку решение не начато.

Первый неполный шаг --- не определена общая область допустимых значений исходного произведения и не рассмотрены его множители.

\textbf{Почему это существенно}

Квадратный корень требует
\[
\cos x\ge 0.
\]
Тангенс требует
\[
\cos x\ne 0.
\]
Вместе эти условия дают
\[
\cos x>0.
\]

Поэтому значения, при которых \(\sqrt{\cos x}=0\), сами по себе не могут быть корнями исходного уравнения: при них тангенс не определён.

\textbf{\idea} для произведения с радикалом и тангенсом сначала объединяем их ограничения в одно условие \(\cos x>0\), а затем пересекаем с ним решения каждого множителя.

\textbf{Исправление}

Начинаем с условия
\[
\cos x>0.
\]
После этого отдельно рассматриваем равенство нулю каждого из трёх множителей и проверяем найденные серии по этому условию.

\textbf{Полное правильное решение}

Исходное уравнение:
\[
(\cos 2x-1)(\operatorname{tg}x-1)\sqrt{\cos x}=0.
\]

Область допустимых значений:
\[
\cos x\ge 0
\]
из квадратного корня и
\[
\cos x\ne 0
\]
из тангенса. Значит,
\[
\cos x>0.
\]

\textbf{1. Первый множитель}
\[
\cos 2x-1=0,
\]
\[
\cos 2x=1.
\]
Отсюда
\[
2x=2\pi n,\qquad n\in\mathbb Z,
\]
\[
x=\pi n.
\]
Теперь проверяем \(\cos x>0\). Для \(x=\pi n\):
\[
\cos(\pi n)=(-1)^n.
\]
Положительное значение получается только при чётных \(n\). Поэтому
\[
x=2\pi n,\qquad n\in\mathbb Z.
\]

\textbf{2. Второй множитель}
\[
\operatorname{tg}x-1=0,
\]
\[
\operatorname{tg}x=1.
\]
Общее решение:
\[
x=\frac{\pi}{4}+\pi n,\qquad n\in\mathbb Z.
\]
Нужно оставить только точки, где \(\cos x>0\). При добавлении \(\pi\) знак косинуса меняется, поэтому допустимы только чётные значения параметра:
\[
x=\frac{\pi}{4}+2\pi n,\qquad n\in\mathbb Z.
\]

\textbf{3. Третий множитель}
\[
\sqrt{\cos x}=0,
\]
\[
\cos x=0.
\]
Но эти значения противоречат условию \(\cos x>0\) и одновременно делают тангенс неопределённым. Поэтому эта ветвь не даёт корней.

\textbf{Проверка}

При
\[
x=2\pi n
\]
имеем \(\cos x=1>0\), первый множитель равен нулю, а тангенс определён.

При
\[
x=\frac{\pi}{4}+2\pi n
\]
имеем \(\cos x>0\) и \(\operatorname{tg}x=1\), поэтому второй множитель равен нулю.

Точки с \(\cos x=0\) в исходное выражение не подставляются, поскольку тангенс там не существует.

\textbf{\result}
\[
\boxed{
x=2\pi n
\quad\text{или}\quad
x=\frac{\pi}{4}+2\pi n,\qquad n\in\mathbb Z.
}
\]

\textbf{Задача на закрепление}

Решите уравнение
\[
(\cos 2x+1)(\operatorname{tg}x+1)\sqrt{\cos x}=0.
\]

\end{document}
